\vspace{-0.7em}
\section{Experiments \& Results}
\vspace{-0.5em}
\label{sec:experiments}

% \vspace{-0.5em}
\subsection{Experimental Setup}
\vspace{-0.5em}
\label{sec:experimental_setup}

\noindent\textbf{Models.} We evaluate RI on 16 models across different families, sizes, and architectures to ensure comprehensive coverage. 
Our open-source models include \emph{Gemma-3-12B}~\citep{gemma3_2025}, \emph{Qwen3-32B/235B}~\citep{qwen3_2025} in both think and no-think modes, \emph{Qwen2.5-72B-Instruct}~\citep{qwen2_5_2024}, \emph{Llama 3.1 70B} \citep{llama3_2024}, \emph{Mistral-Large-Instruct-2411}~\citep{mistral_large_2_2024}, \emph{GLM-4.5 and GLM-4.5-Air}\citep{glm45_2025} and \emph{DeepSeek-V3-0324}~\citep{deepseek_v3_2024}. 
Our proprietary models include \emph{Claude 3.5 haiku} \citep{claude35_2024}, \emph{Claude Sonnet 4} \citep{claude4_2025}, \emph{GPT4.1} and \emph{GPT4.1 mini} \citep{gpt41_2025} and \emph{Gemini 2.5 Flash} and \emph{Gemini 2.5 Flash Lite} \citep{gemini25_2025}. 
We use temperature=0.7 and top-p=0.95 across all models. 
More implementation details are provided in \Secref{sec:detailed_experimental_setup}.

\noindent\textbf{Datasets.} We evaluate RI on three scenarios that require model to make accurate, knowledge-aware refusals: factual question answering, extrinsic hallucination detection (hallucination from training data), and intrinsic hallucination detection (hallucination from context).
(1) We use factual question answering to test models' ability to refuse unknown facts. 
Specifically, we use SimpleQA \citep{wei2024measuringshortformfactualitylarge}, which contains verifiable, atomic factual questions that challenge even frontier LLMs. 
(2) We use extrinsic hallucination detection to test whether models correctly refuse to answer when they cannot recall knowledge from training data. 
For this scenario, we use PreciseWikiQA \citep{bangHalluLensLLMHallucination2025}, a dynamically generated question-answering dataset from Wikipedia snippets. 
PreciseWikiQA tests whether models hallucinate information from their training data, assuming Wikipedia knowledge was included during training. 
We follow \citet{bangHalluLensLLMHallucination2025} to generate 2000 questions for evaluation. 
(3) We use intrinsic hallucination detection to test whether models can faithfully recall information with noisy context. 
For this scenario, we use the 3 datasets from FaithEval \citep{ming2025faithevallanguagemodelstay}. 
However, because the \texttt{Unanswerable} and \texttt{Inconsistency} subsets lack ground truth required for RI computation, we create a 1:1 mixed dataset of PreciseWikiQA and FaithEval to report RI.

\noindent\textbf{Baseline Metrics.} We compare RI against five established metrics for measuring knowledge-aware refusal (Table~\ref{tab:baseline_metrics}): Correct Answer Rate, Refusal Rate, Correct given Attempted (C/A), F-score, and Weighted Score. 
We pick $p=0.2$ for the Weighted Score to balance the accuracy and refusal rate. 
We classify all model outputs into three categories following SimpleQA: (1) Correct, (2) Incorrect, or (3) Not Attempted (refusal).




\noindent\textbf{Adjusting Refusal Rates.} We test RI's consistency by measuring how it changes when models exhibit different refusal rates. 
To this end, we use different system prompts to instruct models to be more conservative or active in answering questions. 
These prompts modify refusal tendencies without degrading the quality of refusal decisions, as shown in \Secref{sec:prompt_design_evaluation}.
Specifically, we use four different system prompts to evaluate each model with varying refusal rates in the first pass, while keep one default prompt that instructs models to answer all questions in the second pass. 
The complete system prompts are provided in \Secref{app:system_prompts}.
